\scanpage{20}
i.e.
\[
 H^i(K^\bullet)\otimes_A M\unsure{\simeq|\overset{\sim}{\longrightarrow}} H^i(K\otimes_A M)\qquad\text{si }i\leq N.
\]

Pour ceci, quitte à tronquer $K^\bullet$ en laissant tomber les $K^j$ pour $j\geq N+2$, on peut supposer $K^\bullet$ borné. Il suffit qu'on ait
\[
 \left\{\begin{aligned}
 \operatorname{Tor}^A_i(K^j,M)&=0 &&\text{pour }i>0,\ \text{tout }j\\
 \operatorname{Tor}^A_i(H^j(K^\bullet),M)&=0 &&\text{pour }i>0,\ \text{tout }j.
 \end{aligned}\right.
\]
Comme les modules $K^j,H^j(K)$ sont de t.f. et en nb fini, on est ramené à \deleted{\unsure{voir que|prouver le}} \inserted{prouver le}

\uline{Lemme} 16 \inserted{Soit $R$ cohérent sur $X$.} $\exists\,U$ dense dans $S=\operatorname{Spec}A$ tel que $p\in U$, $f_1\unsure{--|\ldots} f_{\unsure{n|r}}$ \deleted{suite de \unsure{paramètres}} suite $A_p$-régulière implique
\[
 \operatorname{Tor}^A_i(F,A_p/\underline f A_p)=0\qquad\text{pour }i>0.
\]

Or, pour qu'on ait ceci, il suffit, d'après l'argument \unsure{immédiat} du lemme 6, de savoir que $\operatorname{gr}_J(\unsure{\mathcal O_X|B_X})$ \deleted{est} \inserted{et $\operatorname{gr}_J(A)$ soient} plats sur $A/J$ en les points de $U$ ($J=\text{nilradical de } A$). Or on peut choisir $U$ ainsi, par le \uline{th. de platitude générique}.

Le lemme \deleted{\unsure{14|15}}15 montre que \deleted{quitte à \unsure{en bien rappelant|en bien remplaçant} $\unsure{\Omega|B}$, on peut supposer que l'on} \inserted{l'on est} ramené, pour traiter le cas de $\deleted{\unsure{R^j f}}\unsure{R^j|\mathbb R^j|R^q}f_*(\underline{\operatorname{Ext}}^i(F,G\otimes\Omega))$, \deleted{en changeant} à prouver le
% archive sheet number: 19
